\documentclass[11pt,a4paper]{article}
\usepackage[margin=24mm]{geometry}
\usepackage{amsmath,amssymb,amsthm,booktabs,longtable}
\newtheorem{lemma}{Lemma}\newtheorem{proposition}[lemma]{Proposition}
\newcommand{\Aut}{\operatorname{Aut}}\newcommand{\Out}{\operatorname{Out}}
\title{Odd generators and three O'Nan map counts}
\author{Human Compression project}\date{}
\begin{document}\maketitle
For a finite group $L$, let $x_L(p,q)$ count ordered pairs $(r,s)$
with $|r|=p$, $|s|=q$ and $|rs|=2$, and let $e_L(p,q)$ count those
pairs which generate $L$. We prove the entries
\[
 (3,15):6,\qquad (5,15):340,\qquad (15,15):2834
\]
in the fixed Table~1 of \cite{CL}. The argument uses the previously
acquired character evaluations and complete 581-class subgroup roster
for $G=O'N$. These are borrowed inputs, not classifications proved here.
The numerical entry for $(3,15)$ does not prove its asserted division
into four inner and two outer reflexible cases.

\begin{lemma}\label{odd}
Suppose $L=\langle r,s\rangle$, both $r$ and $s$ have odd order,
and $|rs|=2$. Then $L$ has even order, no subgroup of index two,
and at most one subgroup of index three.
\end{lemma}
\begin{proof}
The first assertion follows from the involution $rs$.
An index-two subgroup is normal, and both odd-order generators
map trivially into its quotient; hence it cannot exist.
For a subgroup $K$ of index three, the coset action has image in
$S_3$. Both generators have image in $A_3$, so transitivity implies
that the image is $A_3$. The point stabilizer is trivial, whence
$K$ is the kernel and is normal. If two such subgroups $K_1,K_2$
were distinct, the map into $L/K_1\times L/K_2$ would have image
$C_3\times C_3$: its two coordinate kernels are distinct and both
coordinate maps are surjective. In this quotient $rs$ is trivial,
so the images of $r$ and $s$ are inverse and generate a cyclic group.
This contradicts surjectivity.
\end{proof}

Every subgroup containing an element of order 15 contains a cyclic
subgroup in roster class 532. The fixed roster has only one class
of order 15; every group of that order is cyclic. Following the
complete maximal-subgroup relation upwards gives exactly 38 classes,
including $G$. The explicit list and the exclusions from Lemma~\ref{odd}
are given below. This uses only reachability; no cover multiplicities
are multiplied along chains.

After these exclusions, the possible proper generated classes are
\[
 6,\ 15,\ 16,\ 87,\ 231,\ 232.                 \tag{1}
\]
The first three are maximal and are respectively $J_1$ and the two
$L_2(31)$ classes. The only exclusion beyond odd order and index two
is for rows 31 and 141. Each has four distinct maximal subgroups of
index three, as the literal multiplicities in the roster show.

The subgroup structures needed for the last three rows can be obtained
from smaller named inputs. Row 87 has exactly one subgroup of order
360, in the named class 167, $A_6$. Each of rows 231 and 232 has
exactly one subgroup of order 60, in the named class 371, $A_5$.
These subgroups are normal, since their conjugates would otherwise
give further subgroups of the same order. We use the classical facts
that $A_5,A_6$ are centerless and perfect and that their outer
automorphism groups have orders 2 and 4. The latter $A_6$ fact is
the same classical automorphism premise already used in the retained
$\mathrm{PSL}_2(9)$ treatment. Thus an extension of either of these
groups by $C_3$ induces only inner automorphisms. If $N$ is the
normal simple subgroup, it follows that
$L=N C_L(N)$, with $N\cap C_L(N)=1$ and $|C_L(N)|=3$. Consequently
\[
 H_{87}\cong A_6\times C_3,\qquad
 H_{231}\cong H_{232}\cong A_5\times C_3.       \tag{2}
\]
In particular, the product descriptions are justified from unique
normal simple subgroups and the stated classical premise, rather
than inferred from equality of orders.

\begin{lemma}\label{extension}
For $N=A_5$ or $A_6$ and $H=N\times C_3$,
\[
 e_H(5,15)=0,\qquad e_H(3,15)=2e_N(3,5),\qquad
 e_H(15,15)=2e_N(5,5).
\]
\end{lemma}
\begin{proof}
If the first generator has order five, it maps trivially to $C_3$.
The product relation forces the second image to be trivial too,
preventing generation. In the other two cases, an element of order
15 must have an order-five component in $N$ and a nonidentity
component in $C_3$; neither $A_5$ nor $A_6$ has an element of order
15. The two $C_3$ components are inverse, giving exactly two choices.
For type $(3,15)$, the first $N$ component has order three, since
an identity component would make the product have order five.
Generation of $H$ requires generation of $N$ by the two projections.
Conversely, if those projections generate $N$, the commutator
subgroup of the generated subgroup is contained in $N\times1$ and
surjects onto the perfect group $N$. It is therefore $N\times1$.
The nonidentity $C_3$ projection then gives the whole direct product.
\end{proof}

The remaining finite calculation is entirely on permutations of
five and six letters. Enumerate all even permutations, test each
ordered pair of the indicated orders, and retain it if its product
has order two. Closure under the two permutations determines the
generated subgroup. The complete literal lists give
\[
\begin{array}{c|rr|rr}
 &x_N(3,5)&e_N(3,5)&x_N(5,5)&e_N(5,5)\\\hline
 A_5&120&120&120&120\\
 A_6&1440&0&2880&1440
\end{array}                                                   \tag{3}
\]
For $A_6$, all 1440 pairs in the first column generate a group of
order 60; in the second type, 1440 generate a group of order 60
and 1440 generate $A_6$. A separate reader reconstructs the even
permutations by cycle parity and recomputes every pair and every
closure. Its composition convention is opposite to that of the
acquisition, which preserves both product order and generated subgroup.
No O'Nan permutation enumeration enters this calculation.

To finish, neither 180 nor 1080 divides the order of $J_1$ or
$L_2(31)$. Hence a pair of one of the three types inside those
maximals cannot generate any other group in (1). Its generated
subgroup is the maximal itself, so its intrinsic count equals its
acquired character count. Every nongenerating pair of $G$ belongs
to exactly one generated subgroup. If $c_i$ denotes the global
conjugacy-class length of row $i$, the complete partition is
\[
 e_G(p,q)=x_G(p,q)-\sum_{i\in\{6,15,16,87,231,232\}}
                         c_i e_{H_i}(p,q).                    \tag{4}
\]
This formula does not count maximal overgroups and requires no
unknown incidence multiplicities. Equations (2)--(3) determine all
its small intrinsic inputs. The arithmetic is displayed as numbers
of pairs divided by $|\Aut(G)|=921631011840$:
\[\begin{array}{c|rrr}
(p,q)&x_G/|\mathrm{Aut}(G)|&\sum c_i e_{H_i}/|\mathrm{Aut}(G)|&e_G/|\mathrm{Aut}(G)|\\\hline
(3,15)&\frac{164}{9}&\frac{110}{9}&6\\
(5,15)&361&21&340\\
(15,15)&\frac{25924}{9}&\frac{418}{9}&2834\\
\end{array}\]

Automorphisms act freely on generating pairs, so the last column
counts the required ordered-pair orbits. There is no further
division for duality, including when $p=q$. All three values agree
with the fixed source. These are additional exact determinations,
not new chirality or principal-refutation assertions.

\paragraph{Attribution and scope.}
The complete subgroup roster, the named $A_5,A_6,J_1,L_2(31)$
inputs, the classical alternating-group automorphism facts, and
the previously accepted character values are explicit borrowed
premises. The contribution here is their use through the odd-generator
obstruction, the central-extension reduction, and the disjoint
generated-subgroup count, with the small literal calculation (3).
No priority claim is made for these elementary general arguments.

\paragraph{Complete candidate list.}
The table gives every proper ancestor of cyclic class 532. An
index-two witness is a subgroup row. For rows 31 and 141, the
displayed multiplicities specify distinct index-three subgroups.
\begin{longtable}{rrl}
Row&Order&Disposition\\\hline\endhead
6&175560&Retained in (1)\\
11&25920&Index two: row 17\\
15&14880&Retained in (1)\\
16&14880&Retained in (1)\\
17&12960&Index two: row 25\\
18&12960&Index two: row 25\\
19&12960&Index two: row 25\\
25&6480&Index two: row 31\\
31&3240&Index three: 87 (4)\\
47&2160&Index two: row 86\\
48&2160&Index two: row 87\\
86&1080&Index two: row 141\\
87&1080&Retained in (1)\\
101&720&Index two: row 164\\
141&540&Index three: 231 (2), 232 (2)\\
144&465&Odd order\\
161&360&Index two: row 233\\
162&360&Index two: row 234\\
164&360&Index two: row 235\\
165&360&Index two: row 231\\
166&360&Index two: row 232\\
230&180&Index two: row 305\\
231&180&Retained in (1)\\
232&180&Retained in (1)\\
233&180&Index two: row 306\\
234&180&Index two: row 306\\
235&180&Index two: row 307\\
236&180&Index two: row 307\\
305&90&Index two: row 396\\
306&90&Index two: row 396\\
307&90&Index two: row 396\\
366&60&Index two: row 465\\
396&45&Odd order\\
463&30&Index two: row 532\\
464&30&Index two: row 532\\
465&30&Index two: row 532\\
532&15&Odd order\\
\end{longtable}


\begin{thebibliography}{9}
\bibitem{CL} T. Connor and D. Leemans, \emph{Algorithmic enumeration
of regular maps}, Ars Math. Contemp. 10 (2016), 211--222.
\bibitem{roster} T. Connor and D. Leemans, \emph{The subgroup lattice
of $O'N$}, fixed 23-page roster. Exact native pages and edition are
bound with the certificate.
\end{thebibliography}
\end{document}
